\documentclass[12pt, titlepage]{article}
\usepackage{amsmath}
\usepackage{amsfonts}
\usepackage{amssymb}
\usepackage{fancyhdr}
\usepackage{cancel}
\usepackage{xcolor}
\usepackage[bidi=basic]{babel}
\babelprovide[main,import]{hebrew}
\babelfont[hebrew]{rm}{Hadasim CLM}
\babelfont[hebrew]{sf}{Miriam Mono CLM}
\usepackage{titlesec}
\newcommand\Ccancel[2][black]{\renewcommand\CancelColor{\color{#1}}\cancel{#2}}
\title{חשבון אינפי 2 - 01040281 \\ גליון 4}
\author{יונתן אבידור - 214269565}
\titleformat{\section}{\Large\sffamily\bfseries}{\thesection}{1em}{}
\begin{document}
\maketitle
\pagestyle{fancy}
\fancyhead{}
\fancyfoot[L]{\text{214269565}}
\part*{שאלה 1}
חשבו את הערך של האינטגרל המוכלל הבא:
\[
	\int^{\infty}_1 \frac{\arctan x}{x^2} \, dx
\]
\part*{פתרון 1}
נשתמש באינטגרציה בחלקים
\[
	\begin{array}{cc}
		u = \arctan x        & v' =  x^{-2}     \\
		u' = \frac{1}{1+x^2} & v = \frac{-1}{x}
	\end{array}
\]
ונקבל
\begin{gather*}
	\lim_{c \to \infty}\int^{c}_{1} \frac{\arctan x}{x^2} \, dx = -\frac{\arctan x}{x}\bigg{|}^c_1 - \int_1^c \frac{-1}{x(1+x^2)} \, dx \\= \lim_{c\to\infty} -\frac{\arctan c}{c} + \underbrace{\arctan{1}}_{\frac{\pi}{4}} + \int_1^c \frac{1}{x(1+x^2)}\,dx\\
	\underset{\text{ארימתמטיקה של גבולות, חיבור}}{=} \lim_{c\to\infty} -\frac{\arctan c}{c} + \frac{\pi}{4} + \lim_{c\to\infty}\int_1^c \frac{1}{x(1+x^2)}\,dx
\end{gather*}
יש לציין שההשלב האחרון יעבוד רק אם אכן נמצא שהאיטנגרל המוכלל מימין קיים, נבדוק זאת:
\begin{gather*}
	\lim_{c\to\infty}\int^c_1 \frac{1}{x\left(1+x^2\right)}\,dx
\end{gather*}
נעשה פירוק לשברים חלקיים
\[
	\frac{1}{x\left(1+x^2\right)} = \frac{A}{x} + \frac{Bx+C}{x^2+1}
\]
נפתח ונקבל
\begin{gather*}
	(A+B)x^2+Cx +A = 0x^2 + 0x + 1
\end{gather*}
לכן $A+B = 0, C= 0, A=1$, מכאן ש-$B=-1$, ומתקבל
\[
	\frac{1}{x(1+x^2)} = \frac{1}{x} - \frac{x}{x^2+1}
\]
ולכן
\begin{gather*}
	\lim_{c\to\infty} \int^{c}_1 \frac{1}{x(1+x^2)}\,dx = \lim_{c\to\infty} \int^{c}_1 \frac{1}{x}-\frac{x}{x^2+1} \\
	\underset{\text{לינאריות, חשבון גבולות} }{=} \lim_{c\to\infty}\int_1^c  \frac{1}{x}\,dx -\lim_{c\to\infty}\int_1^c  \frac{x}{x^2+1}\,dx\\
	\underset{\text{ניוטון לייבניץ}}{=} \lim_{c\to\infty}\log x \bigg{|}^c_1 - \lim_{c\to\infty}\int_1^c  \frac{x}{x^2+1}\,dx
\end{gather*}
נטפל באינטגרל שנותר
\[
	\lim_{c\to\infty} \frac{x}{x^2+1}\,dx \underset{\text{הכפלה וחילוק ב-$2$, לינאריות}}{=} \lim_{c\to\infty} \frac{1}{2} \int^c_1 \frac{2x}{x^2+1}\,dx
\]
נשים לב ש-$\left(x^2+1\right)' = 2x$ ולכן יש לנו נגזרת חלקי הפונקציה ולכן
\[
	\lim_{c\to\infty} \frac{1}{2} \int^c_1 \frac{2x}{x^2+1}\,dx = \lim_{c\to\infty} \frac{1}{2} \log(x^2+1)\bigg{|}^c_1
\]
נשים לב ש-$x^2+1$ חיובית תמיד לכן אין צורך בערך מוחלט\\
נחזיר את הגבולות לביטוי אחד (מאריתמטיקת גבולות, חיבור)
\begin{gather*}
	\lim_{c\to\infty} \log x - \frac{1}{2}\log (x^2+1)\bigg{|}^c_1 \underset{\text{חוקי לוגים}}{=}\lim_{c\to\infty} \log\left(\frac{x}{\sqrt{x^2+1}}\right)\bigg{|}^c_1 \\
	= \lim_{c\to\infty} \log\left(\frac{c}{\sqrt{c^2+1}}\right) - \log\left(\frac{1}{\sqrt{2}}\right)
\end{gather*}
נמצא את הגבול בתוך הלוגריתם
\[
	\lim_{c\to\infty} \frac{c}{\sqrt{c^2+1}} = \lim_{c\to\infty} \frac{\cancel{c}}{\cancel{c}\sqrt{1+\underbrace{\frac{1}{c^2}}_{\to 0}}} = \lim_{c\to\infty} \frac{1}{\sqrt{1}} = 1
\]
לכן
\[
	\lim_{c\to\infty} \log\left(\frac{c}{\sqrt{c^2+1}}\right) - \log\left(\frac{1}{\sqrt{2}}\right) = \log(1) - \log(\frac{1}{\sqrt{2}}) = \frac{\log(2)}{2}
\]
הגבול אכן קיים ולכן החיבור עובד. סה"כ קיבלנו
\[
	\boxed{\int_1^\infty \frac{\arctan x}{x^2} =\frac{\pi}{4} + \frac{\log(2)}{2}}
\]
\newpage
\part*{שאלה 2}
מצאו את כל הערכים של $p>0$ ושל $q>0$ בהם האינטגרלים הבאים מתכנסים
\section*{סעיף א'}
\[
	\int_1^\infty \frac{\sin x}{x^p}\,dx
\]
\section*{סעיף ב'}
\[
	\int_0^\infty \frac{1}{x^p+x^q}\,dx
\]
\section*{סעיף ג'}
\[
	\int_0^\infty \frac{1}{x^p\left(1+x\right)^q}\,dx
\]
\part*{פתרון 2}
\section*{סעיף א'}
נטען כי האינטגרל מתכנס לכל $p>0$, נרצה להשתמש במשפט דיריכלה.\\
נגדיר $f(x) = \sin(x), g(x) = \frac{1}{x^p}$\\
$F(x)$ חסומה\\
נחשב את $g'(x)$
\[
	g(x) = \frac{1}{x^{p}} = x^{-p} \implies g'(x) = -p\cdot x^{p+1} = -\frac{p}{x^{p+1}}
\]
נבדוק את $\int^\infty_1 \left|g'(x)\right|\,dx$, נשים לב שלכל $x$ ולכל $p>0$, $\left|-\frac{p}{x^{p+1}}\right| = \frac{p}{x^{p+1}}$
\[
	\lim_{c\to\infty} \int_1^c\frac{p}{x^{p+1}}\,dx = \lim_{c\to\infty} p \int^c_1 x^{-p-1} \,dx= \lim_{c\to\infty}  \left(\frac{-1}{x^{p}}\right)\bigg{|}^c_1 = \lim_{c\to\infty} \underbrace{\frac{-1}{c^p}}_{\to 0} - \underbrace{\frac{-1}{1^p}}_{\to -1} = 1
\]
הגבול קיים ולכן האינטגרל המוכלל מתכנס.\\
לכל $p>0$ $\lim_{x\to\infty}\frac{1}{x^{p}}=0$ כגבול ידוע מאינפי $1$\\
כל התנאים של משפט דיריכלה מתקיימים, לכן לפי משפט דיריכלה, האינטגרל המוכלל $\int^\infty_1(f\cdot g)(x)=\int^\infty_1 \frac{\sin x}{x^p}\,dx$ מתכנס לכל $\boxed{p>0}$
\section*{סעיף ב'}
הפונקציה "בעייתית" בשני איזורים, כאשר $x\to 0$ וכאשר $x\to \infty$, נרצה לטפל במקרים הללו בנפרד, לכן נפצל את האינטגרל בדרך הבאה:
\[
	\int_0^\infty \frac{1}{x^p+x^q}\,dx \underset{\text{אדטיביות, אם האינטגרלים מתכנסים}}{=}	\int_0^1 \frac{1}{x^p+x^q}\,dx + 	\int_1^\infty \frac{1}{x^p+x^q}\,dx
\]
נשים לב שהפונקציה חיובית בכל תחום האינטגרציה, נשתמש במבחן ההשוואה הגבולי
נגדיר $m:= \min\{p,q\}$
נתבונן ב-$\lim_{x\to 0^{+}} \frac{\frac{1}{x^p+x^q}}{\frac{1}{x^m}}$
\[
	x^{p-m} + x^{q-m}	\lim_{x\to0^{+}} \frac{\frac{1}{x^p+x^q}}{\frac{1}{x^m}} =\lim_{x\to0^{+}} \frac{x^m}{x^p+x^q} = \lim_{x\to0^+}\frac{1}{x^{p-m} + x^{q-m}}
\]
מכיוון ש-$m$ הוא המינימום של $p,q$, אחד מהמעריכים הוא $0$ והשני הוא אי-שלילי, לכן הגבול הוא מספר חיובי כלשהו, לכן מבחן ההשוואה הגבולי תקף.\\
$\frac{1}{x^m}$ מתכנס לכל $m>1$, לכן $\int_0^1 \frac{1}{x^p+x^q}\,dx$ מתכנס אמ"מ $\min\{p,q\}< 1$\\
עבור $\left[1,\infty\right)$ נשתמש גם במבחן ההשוואה הגבולי\\
נגדיר $M:=\max\left\{p,q\right\}$\\
נתבונן ב-$\lim_{x\to \infty} \frac{\frac{1}{x^p+x^q}}{\frac{1}{x^M}}$
\[
	\lim_{x\to \infty} \frac{\frac{1}{x^p+x^q}}{\frac{1}{x^M}} = \lim_{x\to\infty}\frac{x^M}{x^p+x^q} = \lim_{x\to\infty} \frac{1}{x^{p-M}+x^{q-M}}
\]
מכיוון ש-$M$ הוא המקסימום של $p,q$, אחד מהמעריכים הוא $0$ והשני הוא אי-חיובי,\\
לכן הגבול הוא מספר חיובי כלשהו, לכן מבחן ההשוואה הגבולי תקף.\\
$\frac{1}{x^M}$ מתכנס לכל $M>1$, לכן $\int_1^\infty \frac{1}{x^p+x^q}\,dx$ מתכנס אמ"מ $\max\{p,q\} > 1$\\
האינטגרל $\int_0^\infty \frac{1}{x^p+x^q}\,dx$ מתכנס אם ורק אם $\int_0^1 \frac{1}{x^p+x^q}\,dx $ ו-$\int_1^\infty \frac{1}{x^p+x^q}\,dx$ מתכנסים, כלומר כאשר
\[
	\boxed{\min\{p,q\} < 1 < \max\{p,q\}}
\]
\section*{סעיף ג'}
מנימוק זהה לסעיף ב', נפצל את האינטגרל
\[
	\int_0^\infty \frac{1}{x^p\left(1+x\right)^q}\,dx =\int_0^1 \frac{1}{x^p\left(1+x\right)^q}\,dx + \int_1^\infty \frac{1}{x^p\left(1+x\right)^q}\,dx
\]
עבור האינטגרל הראשון נשתמש במבחן ההשוואה הגבולי עם $\frac{1}{x^p}$
$x^p \leq x^p\left(1+x\right)^q$ ולכן $\frac{1}{x^p} \geq \frac{1}{x^p+\left(1+x\right)^q}$
\[
	\lim_{x\to0^+} \frac{\frac{1}{x^p(1+x)^q}}{\frac{1}{x^p}} =\lim_{x\to0^+} \frac{1}{\left(1+x\right)^q} = 1
\]
הגבול קיים וחיובי ולכן ניתן להשתמש במבחן ההשוואה\\
$\int_0^1 \frac{1}{x^p}$ קיים אם ורק אם $p<1$, נעבור לאינטגרל השני\\
נשתמש שוב במבחן ההשוואה הגבולי, אבל הפעם עם $\frac{1}{x^{p+q}}$\\
כאשר $x\geq 1$, $x^{p+q} \leq x^p(1+x)^q$ ולכן $\frac{1}{x^{p+q}} \geq \frac{1}{x^p\left(1+x\right)^q}$
\[
	\lim_{x\to\infty} \frac{\frac{1}{x^p\left(1+x\right)^q}}{\frac{1}{x^{p+q}}} = \lim_{x\to\infty} \frac{x^{p}\cdot x^{q}}{x^p+ x^{p+q}} = \lim_{x\to\infty} \frac{x^q}{\left(1+x\right)^q} =  \lim_{x\to\infty} \left(\frac{x}{1+x}\right)^q = 1
\]
הגבול קיים וחיובי ולכן ניתן להשתמש במבחן ההשוואה\\
$\int_0^1 \frac{1}{x^{p+q}}$ קיים אם ורק אם $p+q > 1$. לכן
\[
	\boxed{p < 1,  q > 1-p}
\]
\newpage
\part*{שאלה 3}
בדקו האם האינטגרלים המוכללים הבאים קיימים
\section*{סעיף א'}
\[
	\int_0^1 \frac{e^x}{x}\,dx
\]
\section*{סעיף ב'}
\[
	\int_0^1 \frac{1}{\log x}\,dx
\]
\section*{סעיף ג'}
\[
	\int_{-\infty}^\infty \frac{2x}{e^x-e^{-x}}\,dx
\]
\part*{פתרון 3}
\section*{סעיף א'}
$\forall x \in \left[0,1 \right]:e^x \geq 1$ ולכן $\frac{e^x}{x} \geq \frac{1}{x}$, בנוסף לכך, שתי הפונקציות אי-שליליות, לכן ניתן להשתמש במבחן ההשוואה.\\
$\int_0^1 \frac{1}{x}\,dx$ מתבדר ולכן $\int_0^1 \frac{e^x}{x}\,dx$ מתבדר גם כן
\section*{סעיף ב'}
יש שתי נקודות בעייתיות. $x=1$ ו-$x=0$ \\
לכן נפצל את האינטגרל
\[
	\int_0^1 \frac{1}{\log x}\,dx = \int_0^\frac{1}{e} \frac{1}{\log x}\,dx  + \int_\frac{1}{e}^1 \frac{1}{\log x}\,dx
\]
נציב $y= \log x$
\begin{gather*}
	y = \log x \\
	e^y = x\\
	dy = \frac{1}{x}\,dx\\
	dx = x \,dy = e^y \,dy\\
	x \to 0^+ \implies y \to -\infty \\
	x = \frac{1}{e} \implies y = -1 \\
	x \to 1^- \implies y \to 0^-
\end{gather*}
ואז
\[
	\int_0^\frac{1}{e} \frac{1}{\log x}\,dx  + \int_\frac{1}{e}^1 \frac{1}{\log x}\,dx = 	\int_{-\infty}^{-1} \frac{e^y}{y}\,dy  + \int_{-1}^0 \frac{e^y}{y} \,dy
\]
לכל $y\in(-\infty,-1]$ $\left|y\right|\geq1$, לכן $\frac{e^y}{\left|y\right|} \leq e^y$, נשתמש במבחן ההשוואה\\
נבדוק אם $\int_{-\infty}^{-1} e^y \, dy$ מתכנס
\[
    \int_{-\infty}^{-1} e^y \, dy  = \lim_{c\to\infty} \int_{-c}^{-1} e^y \, dy \underset{N-L}{=} \lim_{c\to\infty} e^y \bigg{|}^{-1}_{-c}  = \lim_{c\to\infty} e^{-1} - \underbrace{e^{-c}}_{\to 0} = \frac{1}{e}
\]
הגבול קיים לכן $\frac{e^y}{\left|y\right|}$ מתכנס, לכן $\frac{e^y}{y}$ מתכנס בהחלט.\\
אינטגרל שמתכנס בהחלט בהחלט מתכנס ולכן $\frac{e^y}{y}$ מתכנס.\\
נעבור לאינטגרל השני\\
נשתמש במבחן ההשוואה הגבולי
\[
    \lim_{y\to0^-} \frac{\frac{e^y}{y}}{\frac{1}{y}} = \lim_{y\to0^-} e^y = 1
\]
האינטגרל $\int^0_{-1} \frac{1}{y} \, dy$ מתבדר, ולכן ממבחן ההשוואה הגבולי $\frac{e^y}{y}$ מתבדר גם הוא.\\
אינטגרל אחד מתבדר והשני מתכנס, לכן האינטגרל כולו מתבדר.
\section*{סעיף ג'}
נראה שהפונקציה זוגית
\[
    \frac{2(-x)}{e^{-x}-e^{x}} = \frac{-2x}{-(e^x-e^{-x})}= \frac{2x}{e^x-e^{-x}}
\]
הפונקציה זוגית, לכן מספיק להראות כי $\int_0^\infty \frac{2x}{e^x-e^{-x}}\,dx$ קיים\\
נשתמש במבחן ההשוואה הגבולי עם הפונקציה $\frac{1}{x^2}$
\[
    \lim_{x\to\infty}  \frac{\frac{2x}{e^x-e^{-x}}}{\frac{1}{x^2}} = \lim_{x\to\infty} \frac{2x^3}{e^x-\underbrace{e^{-x}}_{\to0}}  = \lim_{x\to\infty} \frac{2x^3}{e^x}
\]
יש פולינום ממעלה שלישית במונה ופונקציה מעריכית במכנה, לכן אם נעשה לופיטל שלוש פעמים נקבל שהגבול הוא $0$\\
לכן ממסקנה של מבחן ההשוואה הגבולי ומזה ש $\int_0^\infty \frac{1}{x^2}\,dx$ מתכנס, $\int^\infty_0 \frac{2x}{e^x-e^{-x}} \, dx$ מתכנס
\newpage
\part*{שאלה 4}
תהא $f$ רציפה ב-$\left[0,1\right]$ ויהיו $0 < a < b$. הראו כי הגבול הבא קיים ומצאו אותו
\[
	\lim_{\varepsilon \to 0^+} \int^{\varepsilon b}_{\varepsilon a} \frac{f\left(x\right)}{x}\,dx
\]
\part*{פתרון 4}
$f$ רציפה מההגדרה, $\frac{1}{x}$ אינטגרבילית בכל $\left[\varepsilon a,\varepsilon b\right]$ וחיובית, לכן ממשפט ערך הביניים האינטגרלי (עם פונקציית משקל כללית), לכל $\varepsilon > 0$ קיים $c_\varepsilon \in \left[\varepsilon a, \varepsilon b\right]$ כך ש
\[
 \lim_{\varepsilon \to 0^+} \int_{\varepsilon a}^{\varepsilon b} \frac{f(x)}{x}  \,dx =\lim_{\varepsilon \to 0^+} f\left(c_\varepsilon\right) \int_{\varepsilon a}^{\varepsilon b} \frac{1}{x}\,dx \underset{\text{אריתמטיקה, כפל}}{=} \lim_{\varepsilon\to0^+} f \left( c_{\varepsilon} \right) \cdot \lim_{\varepsilon \to 0^+} \int_{\varepsilon a}^{\varepsilon b} \frac{1}{x}\,dx
\]
יש לציין שהמעבר האחרון יעבוד אם ורק אם שני הגבולות קיימים\\
כאשר $\varepsilon \to 0^+$
\[
    0\leftarrow \varepsilon a \leq c_\varepsilon \leq \varepsilon b \to 0
\]
לכן, מכלל הסנדוויץ' $\lim_{\varepsilon\to0^+} c_\varepsilon = 0$, ומכאן $\lim_{\varepsilon\to0^+} f\left(c_\varepsilon\right) = f\left(0\right)$
נתבונן בגבול השני
\begin{gather*}
    \lim_{\varepsilon \to 0^+} \int_{\varepsilon a}^{\varepsilon b} \frac{1}{x}\,dx \underset{N-L}{=} \lim_{\varepsilon \to 0^+} \log\left(x\right)\bigg{|}^{\varepsilon a}_{\varepsilon b} =\lim_{\varepsilon \to 0^+} \log(\varepsilon a )- \log(\varepsilon b)  \underset{\left(*\right)}{=}\lim_{\varepsilon \to 0^+} \log\left(\frac{\cancel{\varepsilon} a}{\cancel{\varepsilon} b}\right) \\
        \underset{\left(*\right)}{=} \lim_{\varepsilon \to 0^+} \log\left(\frac{a}{b}\right) \underset{\left(**\right)}{=} \log \left(\frac{a}{b}\right)
\end{gather*}
$\left(*\right)$ - מכיוון ש$\varepsilon \to 0^+$, $\varepsilon \neq 0$ ולכן ניתן לשים $\varepsilon$ במכנה ולצמצם איתו\\
$\left(**\right)$ - קיבלנו ביטוי ללא תלות ב-$\varepsilon$ ולכן אין הבדל בין הביטוי הגבולי לגבול
קיבלנו גבול סופי בשני הביטויים שפיצלנו, לכן האריתמטיקה תקינה, ולכן קיים הגבול.\\
נמצא אותו
\[
\lim_{\varepsilon\to0^+} f \left( c_{\varepsilon} \right) \cdot \lim_{\varepsilon \to 0^+} \int_{\varepsilon a}^{\varepsilon b} \frac{1}{x}\,dx =\boxed{ f\left(0\right) \cdot \log\left(\frac{a}{b}\right)}
\]
\newpage
\part*{שאלה 5}
בדקו האם הטורים הבאים מתכנסים
\section*{סעיף א'}
\[
	\sum_{n=1}^\infty \frac{\left(n!\right)^2}{\left(2n\right)!}
\]
\section*{סעיף ב'}
\[
	\sum_{n=1}^\infty \frac{4n^2+5n}{n\left(n^2+1\right)^{\frac{3}{2}}}
\]
\part*{פתרון 5}
\section*{סעיף א'}
הטור הוא סכום אינסופי של איברים חיוביים ממש, לכן הטור חיובי
נשתמש במבחן המנה
\begin{gather*}
    \lim_{n\to\infty} \frac{\frac{\left(\left(n+1\right)!\right)^2}{\left(2\left(n+1\right)\right)!}}{\frac{\left(n!\right)^2}{\left(2n\right)!}}=
    \lim_{n\to\infty} \frac{\left(\left(n+1\right)!\right)^2\cdot\left(2n\right)!}{\left(2\left(n+1\right)\right)!\left(n!\right)^2}
    \\   =\lim_{n\to\infty} \frac{\left(n!\cdot\left(n+1\right)\right)^2\cdot\left(2n\right)!}{\left(2n+2\right)!\left(n!\right)^2}
    =\lim_{n\to\infty} \frac{\Ccancel[blue]{\left(n!\right)^2}\left(n+1\right)^2\cdot\Ccancel[green]{\left(2n\right)!}}{\left(2n+2\right)\left(2n+1\right)\Ccancel[green]{\left(2n\right)!}\Ccancel[blue]{\left(n!\right)^2}}\\
    = \lim_{n\to\infty} \frac{\cancel{\left(n+1\right)}\left(n+1\right)}{2\cancel{\left(n+1\right)}\left(2n+1\right)}\\
    = \lim_{n\to\infty} \frac{n+1}{4n+2} = \frac{1}{4} < 1
\end{gather*}
הגבול קיים וקטן מ-$1$, לכן ממבחן המנה הטור מתכנס
\section*{סעיף ב'}
גם כאן הטור אי-שלילי, נרצה להשתמש במבחן הההשוואה
\[
    \frac{4n^2+5n}{n\left(n^2+1\right)^{\frac{3}{2}}} = \frac{\cancel{n}\left(4n+5\right)}{\cancel{n}\left(n^2+1\right)^\frac{3}{2}} \leq \frac{4n+5n}{\left(n^2+1\right)^\frac{3}{2}}  = \frac{9n}{\left(n^2+1\right)^\frac{3}{2}} \leq \frac{9n}{\left(n^2\right)^\frac{3}{2}} = \frac{9n}{n^3} = \frac{9}{n^2}
\]
הטור $\sum_{n=1}^\infty \frac{9}{n^2}$ מתכנס כי הוא כפולה בקבוע של $\sum_{n=1}^\infty \frac{1}{n^2}$, אי-שלילי, וכפי שהראינו $\frac{9}{n^2} \geq \frac{4n^2+5n}{n\left(n^2+1\right)^{\frac{3}{2}}}$ לכל $n$\\
לכן ממבחן ההשוואה, גם הטור $\sum_{n=1}^{\infty} \frac{4n^2+5n}{n\left(n^2+1\right)^{\frac{3}{2}}}$ מתכנס
\newpage
\part*{שאלה 6}
בדקו האם הטורים הבאים מתכנסים
\section*{סעיף א'}
\[
	\sum_{n=1}^\infty \left(\frac{1}{\sqrt{2n+1}}-\frac{1}{\sqrt{3n+1}}\right)
\]
\section*{סעיף ב'}
\[
	\sum_{n=1}^\infty \frac{\cos\left( n\pi\right)}{n^\frac{1}{4}}
\]
\part*{פתרון 6}
\section*{סעיף א'}
לכל $n\in \mathbb{N}_+$ \\
$3n + 1 > 2n + 1  > 1$, לכן $\sqrt{3n+1} > \sqrt{2n+1}$ ולכן $\frac{1}{\sqrt{3n+1}} < \frac{1}{\sqrt{2n+1}}$\\
לכן $\frac{1}{\sqrt{2n+1}}-\frac{1}{\sqrt{3n+1}}>0$, כלומר הטור הוא טור אי-שלילי וניתן להשתמש במבחן ההשוואה הגבולי\\
נשווה לטור $\frac{1}{\sqrt{n}}$
\begin{gather*}
    \lim_{n\to\infty} \frac{\frac{1}{\sqrt{2n+1}}-\frac{1}{\sqrt{3n+1}}}{\frac{1}{\sqrt{n}}} \\
    = \lim_{n\to\infty} \sqrt{n}\left( \frac{1}{\sqrt{2n+1}}-\frac{1}{\sqrt{3n+1}}\right) \\
    = \lim_{n\to\infty} \sqrt{n}\left( \frac{1}{\sqrt{n}\sqrt{2+\frac{1}{n}}}-\frac{1}{\sqrt{n}\sqrt{3+\frac{1}{n}}}\right) \\
    = \lim_{n\to\infty} \sqrt{n} \left(\frac{1}{\sqrt{n}}\left(\frac{1}{\sqrt{2+\frac{1}{n}}}-\frac{1}{\sqrt{3+\frac{1}{n}}}\right)\right)\\
    = \lim_{n\to\infty} \frac{1}{\sqrt{2+\underbrace{\frac{1}{n}}_{\to 0}}} - \frac{1}{\sqrt{3+\underbrace{\frac{1}{n}}_{\to 0}}} \\
    = \frac{1}{\sqrt{2}} - \frac{1}{\sqrt{3}}
\end{gather*}
הגבול קיים, לכן הטורים מתכנסים ומתבדרים יחדיו\\
הגבול קיים, והטור $\frac{1}{\sqrt{n}}$ הוא טור מסוג $\frac{1}{n^p}$ כאשר $p=0.5 < 1$ ולכן הטור מתבדר\\
לכן גם הטור $\sum_{n=1}^\infty \left(\frac{1}{\sqrt{2n+1}}-\frac{1}{\sqrt{3n+1}}\right)$ מתבדר
\section*{סעיף ב'}
נשים לב שלכל $n$ זוגי, $\cos\left(n\pi\right)=1$ ולכל $n$ אי-זוגי' $\cos\left(n\pi\right)=-1$, לכן
\[
    \sum_{n=1}^\infty \frac{\cos \left(n\pi\right)}{n^{\frac{1}{4}}} = \sum_{n=1}^\infty \frac{\left(-1\right)^{n}}{n^{\frac{1}{4}}} 
\]
נתבונן ב$\left\{\left|\frac{\left(-1\right)^n}{n^{\frac{1}{4}}}\right|\right\}^\infty_{n=1}=\left\{ \frac{1}{n^{\frac{1}{4}}}\right\}^\infty_{n=1}$\\
הסדרה מונוטונית שואפת ל-$0$ ולכן הטור $\sum_{n=1}^\infty \frac{\left(-1\right)^{n}}{n^{\frac{1}{4}}}$ הוא טור לייבניץ ולכן מתכנס.\\
הטור בערך מוחלט הוא מסוג $\frac{1}{n^p}$ כאשר $p=\frac{1}{4} <1$ ולכן הטור מתכנס \textbf{בתנאי}
\newpage
\part*{שאלה 7}
תהא $\left\{a_n\right\}^\infty_{n=1}$ סדרה חיובית. הוכיחו כי
\[
	\sum_{n=1}^\infty a_n < \infty \iff \sum_{n=1}^\infty \frac{a_n}{1+a_n} < \infty
\]
\part*{פתרון 7}
$\implies$\\
נשים לב כי לכל $n$, $a_n > \frac{a_n}{1+a_n} > 0$, לכן ממבחן ההשוואה הם מתכנסים ומתבדרים יחדיו\\
נתון כי $\sum_{n=1}^\infty a_n < \infty$, לכן $\sum_{n=1}^\infty \frac{a_n}{1+a_n} < \infty$\\
$\impliedby$\\
נתון כי $\sum_{n=1}^{\infty} \frac{a_n}{1+a_n}$ מתכנס, לכן $\lim_{n\to\infty}\frac{a_n}{1+a_n}=0$\\
מכאן שקיים $N\in\mathbb{N}$ כך שלכל $n>N$ $\frac{a_n}{1+a_n} < \frac{1}{2}$, נתבונן באי שוויון זה
\begin{gather*}
    \frac{a_n}{1+a_n} < \frac{1}{2} \implies 2a_n < 1+a_n \implies a_n < 1 \\
    \implies a_n + 1 < 2  \implies \frac{a_n}{1+a_n} > \frac{a_n}{2} \implies 2\left(\frac{a_n}{1+a_n}\right) > a_n
\end{gather*}
$a_n$ קטן מכפל בקבוע של $\frac{a_n}{1+a_n}$, ושתי הסדרות חיוביות. בנוסף, נתון כי $\sum_{n=1}^\infty \frac{a_n}{1+a_n} < \infty$, ולכן ממבחן ההשוואה $a_n < \infty$\\
$\blacksquare$
\newpage
\part*{שאלה 8}
\section*{סעיף א'}
תהא $f$ פונקציה חיובית מונוטונית יורדת בקרן $\left[1,\infty\right)$.\\
יהיו $I_n = \int^n_1 f\left(x\right)\,dx$ ו-$S_n = \sum_{k=1}^n f\left(k\right), n = 1,2,\ldots$.\\
נגדיר $u_n = S_n - I_n$.\\
הוכיחו כי הסדרה $\left\{u_n\right\}_{n=1}^{\infty}$ מתכנסת
\section*{סעיף ב'}
השתמשו בסעיף א' על מנת להוכיח כי הגבול
\[
	\lim_{n\to\infty} \sum_{j=1}^n \frac{1}{j} - \log n
\]
קיים.
\part*{פתרון 8}
\section*{סעיף א'}
נטען כי הסדרה $u_n$ מונוטונית יורדת.\\
נראה כי $u_{n+1} - u_n \leq 0$
\begin{gather*}
    u_{n+1} - u_n = \left(S_{n+1} - I_{n+1}\right) - \left(S_{n} - I_n\right) = \left(S_{n+1} - S_n\right) - \left(I_{n+1} - I_n\right)
\end{gather*}
נסתכל על הצד השמאלי של הביטוי
\[
    S_{n+1} - S_n = \sum_{k=1}^{n+1} f\left(k\right) - \sum_{k=1}^{n} f\left(k\right) = f\left(n+1\right)
\]
ובצד הימני
\[
    I_{n+1} - I_n = \int_{1}^{n+1} f\left(x\right)\,dx - \int_1^{n} f\left(x\right)\,dx \underset{\text{אדטיביות האינטגרל}}{=} \int_n^{n+1} f\left(x\right) \,dx 
\]
לכן
\[
u_{n+1} - u_n = f(n+1) - \int_n^{n+1} f\left(x\right) \,dx
\]
מכיוון ש-$f(x)$ מונוטונית יורדת, האינפימום שלה בקטע $[n,n+1]$ מתקבל ב-$n+1$, לכן האינטגרל בקטע $[n,n+1]$ יהיה גדול או שווה למלבן שגובהו $f\left(n+1\right)$ ורוחבו $1$, כלומר
\[
    \int_n^{n+1} f\left(x\right) \,dx  \geq f\left(n+1\right)\cdot 1
\]
לכן
\[
u_{n+1} - u_n = f\left(n+1\right) - \int_n^{n+1} f\left(x\right) \,dx \leq f\left(n+1\right) - f\left(n+1\right) = 0
\]
עכשיו נראה כי $u_n$ חסומה מלרע.\\
מאדטיביות האינטגרל נקבל ש
\[
    \int_1^n f\left(x\right)\,dx = \sum_{k=1}^{n-1} \int_k^{k+1} f(x)\,dx
\]
מכיוון ש-$f\left(x\right)$ מונוטונית יורדת, הסופרימום שלה בקטע $[k,k+1]$ מתקבל ב-$k$, לכן האינטגרל בקטע $[k,k+1]$ יהיה קטן או שווה למלבן שגבוהו $f\left(k\right)$ ורוחבו $1$, כלומר
\[
\int_k^{k+1} f(x)\,dx \leq f\left(k\right)
\]
לכן
\[
    \int_1^n f\left(x\right)\,dx = \sum_{k=1}^{n-1} \int_k^{k+1} f(x)\,dx \leq \sum_{k=1}^{n-1} f(k)
\]
לכן
\[
    u_n = \sum_{k=1}^n f(k) - \int_1^n f(x) \, dx \geq \sum_{k=1}^n f(k) - \sum_{k=1}^{n-1} f(k) = f(n)
\]
לכן $u_n \geq f_n > 0$\\
$u_n$ חסומה ומונוטונית, לכן ממשפט על סדרות חסומות ומונוטוניות היא מתכנסת\\
$\blacksquare$
\newpage
\section*{סעיף ב'}
נסמן
\begin{gather*}
    f: \left[1,\infty\right)\\
    f(x)  =\frac{1}{x}
\end{gather*}
$f(x)$ חיובית ומונוטונית, ונגדיר $\left\{u_n\right\}_{n=1}^{\infty}$ כמו בסעיף הקודם
\begin{gather*}
    u_n = \sum_{k=1}^n f(k) - \int_1^n f(x)\,dx = \sum_{k=1}^n \frac{1}{n} - \log x \bigg{|}^n_1 \\
    =  \sum_{k=1}^n \frac{1}{n} - \left(\log n - \log 1\right) =  \sum_{k=1}^n \frac{1}{n} - \log n
\end{gather*}
לפי סעיף א'
$$
\lim_{n\to\infty} u_n = \sum_{k=1}^n \frac{1}{n} - \log n
$$
קיים, כנדרש\\
$\blacksquare$
\end{document}
